\chapter{The Recital}

\lettrine{T}{he} maliciousness of the cardinal did not leave much for  the ambassador to say; nevertheless, the word <restoration> had  struck the king, who, addressing the comte, upon whom his eyes had been fixed  since his entrance,—<Monsieur,> said he, <will you have  the kindness to give us some details concerning the affairs of England. You  come from that country, you are a Frenchman, and the orders which I see  glittering upon your person announce you to be a man of merit as well as a man  of quality.>

<Monsieur,> said the cardinal, turning towards the queen-mother,  <is an ancient servant of your majesty's, Monsieur le Comte de la  Fère.>

Anne of Austria was as oblivious as a queen whose life had been mingled with  fine and stormy days. She looked at Mazarin, whose evil smile promised her  something disagreeable; then she solicited from Athos, by another look, an  explanation.

<Monsieur,> continued the cardinal, <was a Treville  musketeer, in the service of the late king. Monsieur is well acquainted with  England, whither he has made several voyages at various periods; he is a  subject of the highest merit.>

These words made allusion to all the memories which Anne of Austria trembled to  evoke. England, that was her hatred of Richelieu and her love for Buckingham; a  Treville musketeer, that was the whole Odyssey of the triumphs which had made  the heart of the young woman throb, and of the dangers which had been so near  overturning the throne of the young queen. These words had much power, for they  rendered mute and attentive all the royal personages, who, with very various  sentiments, set about recomposing at the same time the mysteries which the  young had not seen, and which the old had believed to be forever effaced.

<Speak, monsieur,> said Louis XIV, the first to escape from  troubles, suspicions, and remembrances.

<Yes, speak,> added Mazarin, to whom the little malicious thrust  directed against Anne of Austria had restored energy and gayety.

<Sire,> said the comte, <a sort of miracle has changed the  whole destiny of Charles II. That which men, till that time, had been unable to  do, God resolved to accomplish.>

Mazarin coughed while tossing about in his bed.

<King Charles II.,> continued Athos, <left the Hague neither  as a fugitive nor a conqueror, but as an absolute king, who, after a distant  voyage from his kingdom, returns amidst universal benedictions.>

<A great miracle, indeed,> said Mazarin; <for, if the news  was true, King Charles II, who has just returned amidst benedictions, went  away amidst musket-shots.>

The king remained impassible. Philip, younger and more frivolous, could not  repress a smile, which flattered Mazarin as an applause of his pleasantry.

<It is plain,> said the king, <there is a miracle; but God,  who does so much for kings, monsieur le comte, nevertheless employs the hand of  man to bring about the triumph of His designs. To what men does Charles II principally owe his re-establishment?>

<Why,> interrupted Mazarin, without any regard for the king's  pride—<does not your majesty know that it is to M. Monk?>

<I ought to know it,> replied Louis XIV, resolutely; <and  yet I ask my lord ambassador, the causes of the change in this General  Monk?>

<And your majesty touches precisely the question,> replied Athos;  <for without the miracle of which I have had the honour to speak, General  Monk would probably have remained an implacable enemy of Charles II. God willed  that a strange, bold, and ingenious idea should enter into the mind of a  certain man, whilst a devoted and courageous idea took possession of the mind  of another man. The combinations of these two ideas brought about such a change  in the position of M. Monk, that, from an inveterate enemy, he became a friend  to the deposed king.>

<These are exactly the details I asked for,> said the king.  <Who and what are the two men of whom you speak?>

<Two Frenchmen, sire.>

<Indeed! I am glad of that.>

<And the two ideas,> said Mazarin;—<I am more curious  about ideas than about men, for my part.>

<Yes,> murmured the king.

<The second idea, the devoted, reasonable idea—the least important,  sir—was to go and dig up a million in gold, buried by King Charles I at  Newcastle, and to purchase with that gold the adherence of Monk.>

<Oh, oh!> said Mazarin, reanimated by the word million. <But  Newcastle was at the time occupied by Monk.>

<Yes, monsieur le cardinal, and that is why I venture to call the idea  courageous as well as devoted. It was necessary, if Monk refused the offers of  the negotiator, to reinstate King Charles II in possession of this million,  which was to be torn, as it were, from the loyalty and not the loyalism of  General Monk. This was effected in spite of many difficulties: the general  proved to be loyal, and allowed the money to be taken away.>

<It seems to me,> said the timid, thoughtful king, <that  Charles II could not have known of this million whilst he was in Paris.>

<It seems to me,> rejoined the cardinal, maliciously, <that  his majesty the king of Great Britain knew perfectly well of this million, but  that he preferred having two millions to having one.>

<Sire,> said Athos, firmly, <the king of England, whilst in  France, was so poor that he had not even money to take the post; so destitute  of hope that he frequently thought of dying. He was so entirely ignorant of the  existence of the million at Newcastle, that but for a gentleman—one of  your majesty's subjects—the moral depositary of the million, who  revealed the secret to King Charles II, that prince would still be vegetating  in the most cruel forgetfulness.>

<Let us pass on to the strange, bold and ingenious idea,>  interrupted Mazarin, whose sagacity foresaw a check. <What was that  idea?>

<This—M. Monk formed the only obstacle to the re-establishment of  the fallen king. A Frenchman imagined the idea of suppressing this  obstacle.>

<Oh! oh! but he is a scoundrel, that Frenchman,> said Mazarin;  <and the idea is not so ingenious as to prevent its author being tied up  by the neck at the Place de Grève, by decree of the parliament.>

<Your eminence is mistaken,> replied Athos, dryly; <I did not  say that the Frenchman in question had resolved to assassinate M. Monk, but  only to suppress him. The words of the French language have a value which the  gentlemen of France know perfectly. Besides, this is an affair of war; and when  men serve kings against their enemies they are not to be condemned by a  parliament—God is their judge. This French gentleman, then, formed the  idea of gaining possession of the person of Monk, and he executed his  plan.>

The king became animated at the recital of great actions. The king's  younger brother struck the table with his hand, exclaiming, <Ah! that is  fine!>

<He carried off Monk?> said the king. <Why, Monk was in his  camp.>

<And the gentleman was alone, sire.>

<That is marvellous!> said Philip.

<Marvelous, indeed!> cried the king.

<Good! There are the two little lions unchained,> murmured the  cardinal. And with an air of spite, which he did not dissemble: <I am  unacquainted with these details, will you guarantee their authenticity,  monsieur?>

<All the more easily, my lord cardinal, from having seen the  events.>

<You have?>

<Yes, monseigneur.>

The king had involuntarily drawn close to the count, the Duc d'Anjou had  turned sharply round, and pressed Athos on the other side.

<What next? monsieur, what next?> cried they both at the same time.

<Sire, M. Monk, being taken by the Frenchman, was brought to King Charles  II, at the Hague. The king gave back his freedom to Monk, and the grateful  general, in return, gave Charles II the throne of Great Britain, for which so  many valiant men had fought in vain.>

Philip clapped his hands with enthusiasm, Louis XIV, more reflective, turned  towards the Comte de la Fère.

<Is this true,> said he, <in all its details?>

<Absolutely true, sire.>

<That one of my gentlemen knew the secret of the million, and kept  it?>

<Yes, sire.>

<The name of that gentleman?>

<It was your humble servant,> said Athos, simply, and bowing.

A murmur of admiration made the heart of Athos swell with pleasure. He had  reason to be proud, at least. Mazarin, himself, had raised his arms towards  heaven.

<Monsieur,> said the king, <I shall seek and find means to  reward you.> Athos made a movement. <Oh, not for your honesty, to  be paid for that would humiliate you; but I owe you a reward for having  participated in the restoration of my brother, King Charles II.>

<Certainly,> said Mazarin.

<It is the triumph of a good cause which fills the whole house of France  with joy,> said Anne of Austria.

<I continue,> said Louis XIV: <Is it also true that a single  man penetrated to Monk, in his camp, and carried him off?>

<That man had ten auxiliaries, taken from a very inferior rank.>

<And nothing more but them?>

<Nothing more.>

<And he is named?>

<Monsieur d'Artagnan, formerly lieutenant of the musketeers of your  majesty.>

Anne of Austria coloured; Mazarin became yellow with shame; Louis XIV was  deeply thoughtful, and a drop of moisture fell from his pale brow. <What  men!> murmured he. And, involuntarily, he darted a glance at the minister  which would have terrified him, if Mazarin, at the moment, had not concealed  his head under his pillow.

<Monsieur,> said the young Duc d'Anjou, placing his hand,  delicate and white as that of a woman, upon the arm of Athos, <tell that  brave man, I beg you, that Monsieur, brother of the king, will to-morrow drink  his health before five hundred of the best gentlemen of France.> And, on  finishing those words, the young man, perceiving that his enthusiasm had  deranged one of his ruffles, set to work to put it to rights with the greatest  care imaginable.

<Let us resume business, sire,> interrupted Mazarin, who never was  enthusiastic, and who wore no ruffles.

<Yes, monsieur,> replied Louis XIV <Pursue your  communication, monsieur le comte,> added he, turning towards Athos.

Athos immediately began and offered in due form the hand of the Princess  Henrietta Stuart to the young prince, the king's brother. The conference  lasted an hour; after which the doors of the chamber were thrown open to the  courtiers, who resumed their places as if nothing had been kept from them in  the occupations of that evening. Athos then found himself again with Raoul, and  the father and son were able to clasp each other's hands.  